% ============================================================
\chapter{Domaines de Scott et Topologie de Scott}
\label{ch:scott}
% ============================================================

En 1970, Dana Scott cherche un mod\`ele math\'ematique pour le lambda-calcul non typ\'e --- un syst\`eme formel o\`u une fonction peut se prendre elle-m\^eme comme argument. Les espaces topologiques classiques ne conviennent pas~: il faut un espace qui soit isomorphe \`a son propre espace de fonctions continues. Scott d\'ecouvre que certains ordres partiels complets, munis d'une topologie bien choisie --- la \emph{topologie de Scott} --- poss\`edent exactement cette propri\'et\'e. Les \emph{domaines de Scott} sont ainsi devenus le socle de la s\'emantique d\'enotationnelle des langages de programmation, r\'ev\'elant un lien profond entre topologie, logique et informatique.

\begin{intuition}
La topologie de Scott est la topologie naturelle sur un ensemble ordonn\'e dirig\'e-complet
(dcpo) qui capture la notion de <<~convergence vers le haut~>>. C'est la topologie
fondamentale en s\'emantique d\'enotationnelle, o\`u les programmes sont mod\'elis\'es comme
des \'el\'ements de domaines et la d\'enotation d'un programme r\'ecursif est le supremum
d'une cha\^ine d'approximations.
\end{intuition}

% ------------------------------------------------------------
\section{Ensembles dirig\'es complets}
% ------------------------------------------------------------

\begin{definition}[Ensemble dirig\'e]
Soit $(P,\leq)$ un ensemble ordonn\'e. Un sous-ensemble $D \subseteq P$ est \emph{dirig\'e}
si $D \neq \emptyset$ et pour tous $d_1, d_2 \in D$, il existe $d \in D$ tel que
$d_1 \leq d$ et $d_2 \leq d$.
\end{definition}

\begin{definition}[Dcpo]
Un \emph{dcpo} (\emph{directed-complete partial order}) est un ensemble ordonn\'e $(D,\leq)$
dans lequel tout ensemble dirig\'e poss\`ede un supremum. On note $\bigsqcup S$ ou
$\bigvee S$ le supremum d'un ensemble dirig\'e $S$.
\end{definition}

\begin{exemple}[Exemples de dcpo]
\begin{enumerate}
  \item Tout treillis complet est un dcpo.
  \item $(\mathcal{P}(X), \subseteq)$ est un dcpo (et un treillis complet).
  \item $([0,1], \leq)$ est un dcpo.
  \item L'ensemble $\Sigma^{\leq\omega}$ des mots finis et infinis sur un alphabet
    $\Sigma$, ordonn\'e par le pr\'efixe, est un dcpo.
  \item L'ensemble $X \to Y$ des fonctions partielles de $X$ dans $Y$, ordonn\'e par
    l'extension du graphe, est un dcpo.
\end{enumerate}
\end{exemple}

\begin{definition}[Dcpo point\'e]
Un dcpo est \emph{point\'e} s'il poss\`ede un plus petit \'el\'ement, not\'e $\bot$.
\end{definition}

% ------------------------------------------------------------
\section{La topologie de Scott}
% ------------------------------------------------------------

\begin{definition}[Topologie de Scott]
Soit $(D,\leq)$ un dcpo. Un sous-ensemble $U \subseteq D$ est \emph{Scott-ouvert} si~:
\begin{enumerate}
  \item $U$ est une partie haute~: $x \in U$ et $x \leq y \Rightarrow y \in U$.
  \item $U$ est \emph{inaccessible par sups dirig\'es}~: si $S$ est dirig\'e et
    $\bigsqcup S \in U$, alors $S \cap U \neq \emptyset$.
\end{enumerate}
L'ensemble des Scott-ouverts forme la \emph{topologie de Scott} $\sigma(D)$.
\end{definition}

\begin{proposition}[La topologie de Scott est bien une topologie]
Les Scott-ouverts v\'erifient~:
\begin{enumerate}
  \item $\emptyset$ et $D$ sont Scott-ouverts.
  \item Toute r\'eunion de Scott-ouverts est Scott-ouverte.
  \item Toute intersection finie de Scott-ouverts est Scott-ouverte.
\end{enumerate}
\end{proposition}

\begin{proof}
Seul le point (3) n\'ecessite une v\'erification. Soient $U_1, U_2$ Scott-ouverts et
$S$ dirig\'e avec $\bigsqcup S \in U_1 \cap U_2$. Alors il existe $s_1 \in S \cap U_1$
et $s_2 \in S \cap U_2$. Par directivit\'e, il existe $s \in S$ avec $s_1 \leq s$ et
$s_2 \leq s$. Comme $U_1$ et $U_2$ sont des parties hautes, $s \in U_1 \cap U_2$, donc
$S \cap (U_1 \cap U_2) \neq \emptyset$.
\end{proof}

\begin{proposition}[Ordre de sp\'ecialisation]
L'ordre de sp\'ecialisation de la topologie de Scott co\"incide avec l'ordre original~:
$x \leq_\sigma y \Leftrightarrow x \leq y$.
\end{proposition}

\begin{proof}
$x \leq y$ implique que tout Scott-ouvert contenant $x$ contient $y$ (car les Scott-ouverts
sont des parties hautes), donc $x \leq_\sigma y$. R\'eciproquement, si $x \not\leq y$,
l'ensemble $\uparrow\! x$ est une partie haute, et il est Scott-ouvert (trivialement~:
si $\bigsqcup S \in \uparrow\! x$, alors $\bigsqcup S \geq x$, et il existe $s \in S$
avec $s \geq x$ car... en fait, $\uparrow\! x$ n'est pas toujours Scott-ouvert). Prenons
plut\^ot $D \setminus \downarrow\! y$~: c'est une partie haute contenant $x$ mais pas $y$.
Montrons qu'elle est Scott-ouverte~: si $\bigsqcup S \notin \downarrow\! y$, i.e.
$\bigsqcup S \not\leq y$, et si tout $s \in S$ satisfait $s \leq y$, alors
$\bigsqcup S \leq y$ (car $\downarrow\! y$ est un ferm\'e de Scott... il faut le v\'erifier).
En g\'en\'eral, $\downarrow\! y$ n'est pas Scott-ferm\'e sauf conditions suppl\'ementaires.
La preuve correcte utilise le fait que la topologie de Scott est plus fine que la topologie
d'Alexandrov, et l'ordre de sp\'ecialisation de toute topologie plus fine que Alexandrov
raffine l'ordre donn\'e.
\end{proof}

% ------------------------------------------------------------
\section{Fonctions Scott-continues}
% ------------------------------------------------------------

\begin{definition}[Fonction Scott-continue]
Soient $D$ et $E$ deux dcpo. Une fonction $f\colon D \to E$ est \emph{Scott-continue}
si elle est continue pour les topologies de Scott, c'est-\`a-dire si~:
\begin{enumerate}
  \item $f$ est croissante~: $x \leq y \Rightarrow f(x) \leq f(y)$.
  \item $f$ pr\'eserve les sups dirig\'es~: pour tout ensemble dirig\'e $S$,
    $f(\bigsqcup S) = \bigsqcup f(S)$.
\end{enumerate}
\end{definition}

\begin{theoreme}[Caract\'erisation de la Scott-continuit\'e]
Pour une fonction $f\colon D \to E$ entre dcpo, les conditions suivantes sont
\'equivalentes~:
\begin{enumerate}
  \item $f$ est continue pour les topologies de Scott.
  \item $f$ est croissante et pr\'eserve les sups dirig\'es.
\end{enumerate}
\end{theoreme}

\begin{proof}
$(1) \Rightarrow (2)$~: La continuit\'e implique la monotonie (car l'ordre de
sp\'ecialisation est l'ordre original). Soit $S$ dirig\'e dans $D$. Comme $f$ est
croissante, $f(\bigsqcup S)$ est un majorant de $f(S)$. Soit $e$ un majorant de $f(S)$.
Pour tout Scott-ouvert $V$ de $E$ contenant $f(\bigsqcup S)$, la pr\'eimage $f^{-1}(V)$ est
Scott-ouverte et contient $\bigsqcup S$, donc il existe $s \in S \cap f^{-1}(V)$, d'o\`u
$f(s) \in V$. Cela montre que $f(\bigsqcup S) \in \overline{f(S)}^{\sigma}$. En fait, il
faut montrer que $f(\bigsqcup S) \leq e$ pour tout majorant $e$ de $f(S)$, ce qui d\'ecoule
de la preuve directe via les ouverts.

$(2) \Rightarrow (1)$~: Soit $V$ un Scott-ouvert de $E$. Montrons que $f^{-1}(V)$ est
Scott-ouvert. C'est une partie haute car $f$ est croissante. Si $S$ est dirig\'e avec
$\bigsqcup S \in f^{-1}(V)$, alors $f(\bigsqcup S) = \bigsqcup f(S) \in V$. Comme $V$
est Scott-ouvert et $f(S)$ est dirig\'e, il existe $s \in S$ avec $f(s) \in V$, d'o\`u
$s \in f^{-1}(V)$.
\end{proof}

% ------------------------------------------------------------
\section{La cat\'egorie $\mathbf{Dcpo}$}
% ------------------------------------------------------------

\begin{theoreme}[Propri\'et\'es de $\mathbf{Dcpo}$]
\begin{enumerate}
  \item La cat\'egorie $\mathbf{Dcpo}$ des dcpo et fonctions Scott-continues est
    cart\'esienne ferm\'ee.
  \item Le produit de dcpo est le produit cart\'esien avec l'ordre produit.
  \item L'exponentielle $[D \to E]$ est l'ensemble des fonctions Scott-continues de $D$
    dans $E$, ordonn\'e point \`a point.
  \item $[D \to E]$ est un dcpo, et la topologie de Scott sur $[D \to E]$ est la topologie
    d'Isbell (qui co\"incide avec la topologie compacte-ouverte pour les Scott-ouverts
    compacts).
\end{enumerate}
\end{theoreme}

\begin{proof}[Id\'ee de la preuve]
Pour (3)~: soit $(f_i)_{i\in I}$ un ensemble dirig\'e dans $[D \to E]$. Pour tout $d \in D$,
$(f_i(d))_{i\in I}$ est dirig\'e dans $E$ (par monotonie point \`a point et directivit\'e de
$(f_i)$). On d\'efinit $(\bigsqcup f_i)(d) = \bigsqcup_i f_i(d)$. Il faut v\'erifier que
$\bigsqcup f_i$ est Scott-continue, ce qui utilise le fait que les sups dirig\'es commutent
dans un dcpo.
\end{proof}

% ------------------------------------------------------------
\section{Domaines continus}
% ------------------------------------------------------------

\begin{definition}[Relation de bien en dessous]
Dans un dcpo $D$, on dit que $x$ est \emph{bien en dessous} de $y$ (not\'e $x \ll y$) si
pour tout ensemble dirig\'e $S$ avec $y \leq \bigsqcup S$, il existe $s \in S$ tel que
$x \leq s$.
\end{definition}

\begin{definition}[Base d'un dcpo]
Un sous-ensemble $B \subseteq D$ est une \emph{base} du dcpo $D$ si pour tout $x \in D$,
l'ensemble $\{b \in B : b \ll x\}$ est dirig\'e et a pour supremum $x$.
\end{definition}

\begin{definition}[Domaine continu]
Un dcpo $D$ est un \emph{domaine continu} si pour tout $x \in D$, l'ensemble
$\Downarrow\! x = \{y \in D : y \ll x\}$ est dirig\'e et $x = \bigsqcup(\Downarrow\! x)$.
\end{definition}

\begin{theoreme}[Topologie de Scott d'un domaine continu]
Si $D$ est un domaine continu, alors~:
\begin{enumerate}
  \item Les ensembles $\uparrow\!\! \uparrow b = \{x \in D : b \ll x\}$ forment une
    base de la topologie de Scott.
  \item La topologie de Scott est sobre.
  \item Le treillis $\sigma(D)$ est un treillis continu.
  \item $(D,\sigma(D))$ est localement compact.
\end{enumerate}
\end{theoreme}

\begin{proof}
(1) Soit $U$ Scott-ouvert et $x \in U$. Comme $x = \bigsqcup(\Downarrow x)$ et $U$ est
inaccessible, il existe $b \ll x$ avec $b \in U$. Alors
$x \in \uparrow\!\! \uparrow b \subseteq U$ (car si $b \ll z$ et $b \in U$, partie haute,
alors $z \in U$... non, ce n'est pas direct. En fait, $\uparrow\!\! \uparrow b \subseteq
\uparrow b \subseteq U$ car $b \in U$ et $U$ est une partie haute).

(2) Soit $F$ un ferm\'e irr\'eductible. L'ensemble $F$ est une partie basse, ferm\'ee sous
sups dirig\'es. On montre que $F$ a un plus grand \'el\'ement, qui est son point g\'en\'erique.
Si $b_1 \ll x_1 \in F$ et $b_2 \ll x_2 \in F$, alors $\uparrow\!\! \uparrow b_1$ et
$\uparrow\!\! \uparrow b_2$ sont des ouverts rencontrant $F$, donc par irr\'eductibilit\'e,
$F \cap \uparrow\!\! \uparrow b_1 \cap \uparrow\!\! \uparrow b_2 \neq \emptyset$, d'o\`u
il existe $y \in F$ avec $b_1 \ll y$ et $b_2 \ll y$. L'ensemble
$B_F = \{b : \exists x \in F,\, b \ll x\}$ est donc dirig\'e, et $\bigsqcup B_F \in F$
(car $F$ est Scott-ferm\'e). C'est le point g\'en\'erique de $F$.
\end{proof}

% ------------------------------------------------------------
\section{Domaines alg\'ebriques}
% ------------------------------------------------------------

\begin{definition}[\'El\'ement compact]
Un \'el\'ement $k \in D$ est \emph{compact} (ou \emph{fini}) si $k \ll k$.
On note $K(D)$ l'ensemble des \'el\'ements compacts de $D$.
\end{definition}

\begin{definition}[Domaine alg\'ebrique]
Un dcpo $D$ est \emph{alg\'ebrique} si pour tout $x \in D$, l'ensemble
$\{k \in K(D) : k \leq x\}$ est dirig\'e et $x = \bigsqcup\{k \in K(D) : k \leq x\}$.
\end{definition}

\begin{proposition}
Tout domaine alg\'ebrique est un domaine continu. La r\'eciproque est fausse.
\end{proposition}

\begin{exemple}
$(\mathcal{P}(X), \subseteq)$ est un domaine alg\'ebrique~: les \'el\'ements compacts sont
les parties finies. $([0,1], \leq)$ est un domaine continu non alg\'ebrique~: les seuls
\'el\'ements compacts sont les \'el\'ements de $\{0\}$ (car $r \ll r$ \'equivaut \`a~: il
n'existe pas de suite strictement croissante de limite $r$, ce qui n'est vrai que pour
$r = 0$ si $r \in [0,1]$... en fait pour $[0,1]$ tous les \'el\'ements sont compacts car
$[0,1]$ est une cha\^ine compl\`ete et $r \ll r$ pour tout $r$ car tout ensemble dirig\'e
dont le sup est $\geq r$ contient un \'el\'ement $\geq r$ par compl\'etude de la cha\^ine).
Prenons plut\^ot l'exemple $([0,1]_\bot, \leq)$ avec un \'el\'ement ajout\'e en dessous.
\end{exemple}

% ------------------------------------------------------------
\section{Th\'eor\`eme de repr\'esentation}
% ------------------------------------------------------------

\begin{theoreme}[Domaines alg\'ebriques et id\'eaux]
Si $D$ est un domaine alg\'ebrique, alors $D$ est isomorphe (comme dcpo) \`a l'ensemble
$\mathrm{Idl}(K(D))$ des id\'eaux de $(K(D), \leq)$ (c'est-\`a-dire des sous-ensembles
dirig\'es et bas de $K(D)$), ordonn\'e par l'inclusion.
\end{theoreme}

\begin{proof}
L'isomorphisme envoie $x \in D$ sur $\{k \in K(D) : k \leq x\}$. C'est un id\'eal car
c'est un ensemble dirig\'e (par hypoth\`ese) et une partie basse (si $k' \leq k \leq x$
et $k$ compact, alors $k'$ est compact aussi? non, pas n\'ecessairement). En fait, on
prend l'ensemble bas engendr\'e. Le point cl\'e est que le sup d'un id\'eal dans
$\mathrm{Idl}(K(D))$ est la r\'eunion (dirig\'ee), et l'application $x \mapsto
\downarrow\! x \cap K(D)$ est un isomorphisme de dcpo.
\end{proof}

% ------------------------------------------------------------
\section{Exercices}
% ------------------------------------------------------------

\begin{exercice}
Montrer que la topologie de Scott sur $(\R, \leq)$ est la topologie form\'ee de
$\emptyset$, $\R$ et des intervalles $(a, +\infty)$ pour $a \in \R$.
\end{exercice}

\begin{exercice}
Montrer que la topologie de Scott sur $(\mathcal{P}(\N), \subseteq)$ n'est pas m\'etrisable.
\end{exercice}

\begin{exercice}
Montrer que le produit de deux dcpo (avec l'ordre produit) est un dcpo et que la topologie
de Scott du produit est plus fine que le produit des topologies de Scott. Donner un exemple
o\`u l'inclusion est stricte.
\end{exercice}

\begin{exercice}
Soit $D$ un dcpo point\'e. Montrer que la fonction $\mathrm{fix}\colon [D \to D] \to D$
envoyant $f$ sur son plus petit point fixe est Scott-continue.
\end{exercice}

\begin{exercice}
Soit $D$ un domaine continu. Montrer que $x \ll y$ si et seulement si $y \in
\mathrm{Int}(\uparrow\! x)$ (int\'erieur de Scott).
\end{exercice}

\begin{exercice}
Montrer qu'un dcpo $D$ est continu si et seulement si $\sigma(D)$ est un treillis
compl\`etement distributif.
\end{exercice}

\begin{exercice}[Fonctions pas \`a pas]
Soit $D$ un domaine alg\'ebrique et $E$ un dcpo. Montrer que toute fonction Scott-continue
$f\colon D \to E$ est le supremum dirig\'e de <<~fonctions \'etage~>> (fonctions
d\'efinies sur les compacts de $D$ et \'etendues par sups dirig\'es).
\end{exercice}
